\documentclass[11pt]{article}

\usepackage[T1]{fontenc}
\usepackage{amsmath,amssymb,graphicx}
\usepackage[margin=1in]{geometry}
\usepackage{booktabs}
\usepackage{hyperref}

\title{No-Distance Combinatorial Austerity Line from Metric Removal to Harmonic-Address Selection and Locking Null Results on an Edge-Graph Control Branch}
\author{Tomislav Rupi\'c \\ Independent Researcher}
\date{March 2026}

\begin{document}

\maketitle
\begin{center}
\large Paper 25.1
\end{center}

\begin{abstract}
This paper updates the HAOS-IIP combinatorial austerity line from Stage C0.1 through Stage C0.14. The line begins as a strict no-distance control: coordinate distance is removed from the interaction kernel and replaced by graph-internal structure on an explicit edge-graph control branch. It then expands in two directions. The first direction, Stages C0.5--C0.9, asks whether combinatorial survivor rules, incidence-seeded mismatch accumulation, and softer bounded-path kernels can produce sequencing, explicit kernel inconsistency, or broad irreversibility. The second direction, Stages C0.10--C0.14, adds a static harmonic-address compatibility factor on the clustered Dirac--Kahler seed and asks whether topology selection sharpens into a resonance ladder, a detuning phase boundary, or a locking mechanism. The combined result is bounded and mixed. Graph-structural class survival survives metric removal, and incidence-forced protected sequencing appears reproducibly under external combinatorial forcing. A weaker bounded-path survivor kernel can open explicit adjacency-versus-kernel inconsistency in five of nine runs, but that partial positive does not scale under tri-root forcing. The harmonic-address branch shows a weak but real topology-selection signal: exact-match sectors preserve the braid class, resonance alignment stratifies into address bands, and a braid-to-smear transition appears near detuning $\delta \approx 0.75$. But no tested bounded latch converts that selection into locking. The updated boundary is therefore sharp: the C0 line supports graph-structural class survival, limited protected sequencing, and weak harmonic-address topology selection, but it still does not support broad irreversibility scaling or topology locking on the no-distance control branch.
\end{abstract}

\section{Introduction}

The C0 line is a bounded control program inside HAOS-IIP, not a replacement core. Its purpose is narrower:
\begin{quote}
\emph{What survives once coordinate distance is removed, and what further combinatorial ingredients are still insufficient to turn selection into irreversibility or locking?}
\end{quote}

The early answer was reported in Paper 24.1, which froze Stages C0.1 through C0.4 as a metric-removal and topology-control sequence. The present paper extends that line through C0.14 while keeping the core/model boundary explicit:
\begin{itemize}
\item the frozen core remains the HAOS-IIP axioms of recoverability, interaction primacy, representational regimes, viability-constrained dynamics, and invariant residue;
\item the C0 line remains a regime-specific control branch beneath that core;
\item positive and negative numerical outcomes are treated as branch boundaries, not as revisions of the axioms.
\end{itemize}

\section{Core Boundary and Spectral Framing}

The recent canonical cleanup matters here. On Laplacian-style branches, HAOS recoverability can be written as an operator-level spectral condition rather than as free-floating language. If
\[
L = D-A, \qquad 0=\lambda_0 \le \lambda_1 \le \lambda_2 \le \cdots,
\]
then the cleanest scalar-branch compression is
\[
e^{-\lambda_1 T}\delta < \epsilon_c,
\]
where $\delta$ is perturbation size, $T$ is the admissible recovery horizon, and $\epsilon_c$ is the viable coherence threshold. In words: the low end of the interaction spectrum must damp perturbations back below the viable coherence band quickly enough to restore operational structure.

This sharpens the language of recoverability on Laplacian branches, but it does not rewrite any C0 outcome. It clarifies how the scalar control line should be read:
\begin{itemize}
\item C0 positives remain positives only if recoverability improves inside the branch;
\item C0 nulls remain nulls even if a branch looks visually structured;
\item a partial opening of kernel inconsistency is not the same thing as broad irreversibility;
\item topology selection is not the same thing as topology locking.
\end{itemize}

\section{Design Discipline Across the C0 Line}

The C0 line actually contains two sub-branches that share a common discipline but answer different questions.

\subsection{C0.1--C0.9: pure no-distance control branch}

These stages use an explicit projected-transverse edge-graph control operator. The graph size, representative encounter families, and shell-weight measurement vocabulary are fixed while the kernel or survivor rule is varied using only graph-internal quantities:
\begin{itemize}
\item adjacency,
\item bounded path length in edge count,
\item local degree,
\item local motif incidence,
\item combinatorial Laplacian structure.
\end{itemize}

The three representative families are:
\begin{itemize}
\item \texttt{clustered\_composite\_anchor},
\item \texttt{counter\_propagating\_corridor},
\item \texttt{phase\_ordered\_symmetric\_triad}.
\end{itemize}

\subsection{C0.10--C0.14: harmonic-address extension}

These stages shift to the clustered Phase III Dirac--Kahler seed while keeping the projected-transverse sector, graph resolution, and topology measurement stack fixed. A static harmonic-address compatibility factor is added as a relational multiplier on the combinatorial operator. The branch therefore asks a different question:
\begin{quote}
\emph{Can harmonic-address compatibility select, stratify, or lock topology without metric distance?}
\end{quote}

\section{Branch A: Metric Removal and Topology Controls (C0.1--C0.4)}

Paper 24.1 already froze the first four stages. The present synthesis keeps that result intact:
\begin{itemize}
\item C0.1: primary persistence-class behavior survives removal of coordinate-distance weighting on the explicit edge-graph control branch, although transport texture shifts.
\item C0.2: degree-spectrum variation reshapes transport span and flow concentration but does not overturn the coarse topology labels.
\item C0.3: the no-distance branch does not produce a clean braid baseline; the nearest honest surrogate, the counter-propagating corridor, remains \texttt{smeared\_transfer} under all tested incidence-noise perturbations.
\item C0.4: small local motifs do not seed a reproducible new topology class; eight of nine runs are neutral and the only non-neutral response is an isolated diamond-induced channel-splitting case.
\end{itemize}

The stable conclusion from C0.1--C0.4 therefore remains:
\begin{quote}
\emph{The no-distance branch supports graph-structural class survival and transport reshaping, but not a reproducible combinatorial braid family or a motif-seeded topology mechanism.}
\end{quote}

\section{Branch B: Sequencing, Mismatch, and Kernel Inconsistency (C0.5--C0.9)}

\subsection{C0.5: static sequencing baseline}

Stage C0.5 replaces the assumed packet-evolution clock with a combinatorial sequencing rule built from persistence-class stabilization, incidence-mismatch accounting, kernel-update consistency, and protected topology-stability blocks. The result is mostly static. Eight of nine runs freeze into \texttt{static\_regime}, the mismatch ledger stays zero everywhere, and only one protected ordering event appears on the symmetric triad.

This is a useful negative baseline: the unforced branch is too rigid to generate sequencing internally.

\subsection{C0.6: incidence-seeded mismatch accumulation}

Stage C0.6 deliberately injects the noise and motif menus from C0.3 and C0.4 after the protected stability constraint is activated. Here the branch opens clearly:
\begin{itemize}
\item positive mismatch density appears in $9/9$ runs,
\item protected arrow forcing appears in $6/9$ runs,
\item the protected family labels remain intact: three \texttt{trapped\_local}, three \texttt{smeared\_transfer}, three \texttt{split\_channel\_exchange}.
\end{itemize}

This is the strongest positive result in the sequencing line. The branch does support incidence-forced protected ordering. But the mechanism is still not inconsistency-mediated because kernel-consistency failures remain zero.

\subsection{C0.7: original alternation remains too rigid}

Stage C0.7 adds explicit adjacency-versus-graph-shell alternation on top of the positive-mismatch runs from C0.6. The result is cleanly negative:
\begin{itemize}
\item explicit consistency failures remain $0/9$,
\item protected arrow under inconsistency remains $0/9$,
\item the general protected ordering inherited from C0.6 survives only as an ordering-without-inconsistency effect.
\end{itemize}

So the original survivor-kernel pair is too rigid to support explicit inconsistency as the driver.

\subsection{C0.8 and C0.9: softer survivor maps partially open, then fail to scale}

Stage C0.8 replaces the rigid graph-shell survivor map with a bounded-path kernel weighted by $1/(h+1)$ up to a fixed hop cap and upgrades motifs to bi-root and two-layer injections. This is the first stage that opens explicit inconsistency:
\begin{itemize}
\item explicit consistency failures appear in $5/9$ runs,
\item protected-arrow events under weaker inconsistency appear in $2/9$ runs,
\item the hybrid noise-plus-motif protocol is the only forcing family that produces protected ordering under explicit inconsistency, with positive cases on the clustered anchor and the symmetric triad.
\end{itemize}

Stage C0.9 tries to scale that partial opening with tri-root and three-layer motifs. It fails:
\begin{itemize}
\item explicit consistency failures drop to $4/9$,
\item tri-root-specific failures appear in only $2/9$,
\item protected-arrow events under weaker or tri-root inconsistency fall to $0/9$,
\item the irreversibility scaling factor is $0.0$ relative to the C0.8 partial positive.
\end{itemize}

\begin{table}[t]
\centering
\scriptsize
\begin{tabular}{@{}lcccc@{}}
\toprule
Stage & Positive mismatch & Explicit inconsistency & Protected arrow & Main read \\
\midrule
C0.5 & $0/9$ & $0/9$ & $1/9$ & mostly static sequencing baseline \\
C0.6 & $9/9$ & $0/9$ & $6/9$ & incidence-forced protected ordering \\
C0.7 & $9/9$ & $0/9$ & $0/9$ & original alternation too rigid \\
C0.8 & $9/9$ & $5/9$ & $2/9$ & partial inconsistency opening under weaker kernel \\
C0.9 & $9/9$ & $4/9$ & $0/9$ & tri-root scaling fails \\
\bottomrule
\end{tabular}
\caption{Sequencing and inconsistency block. The branch supports reproducible mismatch accumulation and a partial inconsistency opening under a softer bounded-path survivor map, but not broad irreversibility scaling.}
\label{tab:c0b}
\end{table}

The correct combined boundary for C0.5--C0.9 is therefore:
\begin{quote}
\emph{The no-distance branch supports incidence-forced protected sequencing and can be pushed into explicit inconsistency under softer survivor maps, but it does not support a broad claim that irreversibility follows automatically from that inconsistency.}
\end{quote}

\section{Branch C: Harmonic-Address Selection, Detuning, and Locking (C0.10--C0.14)}

\subsection{C0.10: weak proto-selection, no persistence gain}

Stage C0.10 adds a static harmonic-address compatibility factor to the clustered Dirac--Kahler seed. This opens the first clean topology-selection signal:
\begin{itemize}
\item exact-match runs preserve the braid class in $2/9$ runs,
\item near-match and incompatible sectors relax to \texttt{transfer\_smeared},
\item persistence gain remains exactly $0.0$ across the matrix.
\end{itemize}

So harmonic addressing acts as a weak proto-selection rule, not as a persistence generator.

\subsection{C0.11 and C0.12: ladder plus threshold}

Stage C0.11 asks whether discrete harmonic-address differences generate a resonance ladder. The answer is yes, but in a split sense:
\begin{itemize}
\item topology shows two visible bands: same-address versus all nonzero address differences,
\item resonance alignment further stratifies the nonzero sector into $d1/d4$ and $d2/d3$.
\end{itemize}
The same-address braid sector is not fully refinement-stable, so the ladder is stronger in alignment than in topology stability.

Stage C0.12 then scans a denser detuning ladder. The result is a sharp threshold:
\begin{itemize}
\item \texttt{braid\_like\_exchange} persists through $\delta = 0.625$ in $6/9$ runs,
\item the topology flips once near $\delta = 0.750$,
\item the remaining three runs are \texttt{transfer\_smeared} for $\delta = 0.750, 0.875, 1.000$.
\end{itemize}

\subsection{C0.13 and C0.14: boundary refinement and locking null}

Stage C0.13 resolves the detuning corridor near the observed transition. The local picture is sharp:
\begin{itemize}
\item in the baseline, motif-weak-minus, and beta-weak families, \texttt{address\_protected\_braid} survives at $\delta = 0.66, 0.69$ and flips to \texttt{transfer\_smeared} at $\delta = 0.72$,
\item motif-weak-plus, degree-mild, and degree-asymmetric collapse the protected window before the scan even enters the corridor.
\end{itemize}

The global verdict remains \texttt{fragmented multi-window protection} because no family satisfies the strict three-consecutive-sample resolution rule. The boundary is visually sharp in the baseline-like families but not formally resolved by the atlas criterion.

Stage C0.14 then tests the next honest escalation: can bounded local latches convert topology selection into locking? The answer is no.
\begin{itemize}
\item all $9/9$ runs remain \texttt{no\_locking\_effect},
\item no \texttt{locked\_braid} or \texttt{quasi\_bounded\_exchange} class appears,
\item the exact-match family stays \texttt{braid\_like\_exchange},
\item the boundary and fragile families stay \texttt{transfer\_smeared},
\item the harmonic phase latch actually shortens braid survival on the exact-match seed by $0.1253$.
\end{itemize}

\begin{table}[t]
\centering
\scriptsize
\begin{tabular}{@{}lccc@{}}
\toprule
Stage & Main class counts & Strongest positive & Boundary \\
\midrule
C0.10 & $2$ braid, $7$ smear & exact-match address preserves braid & no persistence gain \\
C0.11 & $2$ braid, $7$ dispersive & resonance ladder in alignment & same-address braid not refinement-stable \\
C0.12 & $6$ braid, $3$ smear & sharp flip near $\delta = 0.750$ & not a smooth decay and not multi-plateau \\
C0.13 & $6$ protected braid, $36$ smear & narrow local protected window & no family satisfies strict 3-sample rule \\
C0.14 & $3$ braid, $6$ smear & none & all $9/9$ runs are \texttt{no\_locking\_effect} \\
\bottomrule
\end{tabular}
\caption{Harmonic-address block. Harmonic addressing selects and stratifies topology, but the branch still closes on a locking null result.}
\label{tab:c0c}
\end{table}

The correct combined boundary for C0.10--C0.14 is:
\begin{quote}
\emph{Harmonic addressing produces weak topology selection and a real detuning threshold on the clustered seed, but the tested bounded latch rules do not convert that selection into locking.}
\end{quote}

\section{What Changes and What Does Not}

The updated paper changes the map of the C0 line, but not the core or the main phase boundaries.

\subsection{What changes}

\begin{itemize}
\item The C0 line is no longer just a metric-removal control. It now includes a full sequencing/inconsistency block and a full harmonic-address selection/locking block.
\item The strongest late positive is no longer ``metric removal preserves some classes.'' It is C0.6: incidence-forced protected sequencing with positive mismatch in $9/9$ runs and protected arrow forcing in $6/9$ runs.
\item The strongest late negative is no longer just ``motifs do not seed braid.'' It is C0.14: selection does not become locking under the tested bounded latches.
\end{itemize}

\subsection{What does not change}

\begin{itemize}
\item The frozen HAOS-IIP core axioms do not change.
\item Paper 24.1 remains the frozen early synthesis for C0.1--C0.4.
\item Phase I remains closed by the same boundary: frozen pre-geometry can organize, but it does not bind.
\item Phase II is still not rescued by the C0 line. C0.8 opens only a partial inconsistency signal and C0.9 fails to scale it.
\item Phase III topology findings remain topology findings; the harmonic-address branch adds a selector on that clustered seed but not a persistence law.
\end{itemize}

\section{Conclusion}

The full C0.1--C0.14 line now has a stable architecture and a clean final boundary.

First, metric removal does not erase the no-distance branch. Real graph-structural class survival remains after coordinate distance is removed. Degree landscape, incidence noise, and motif forcing reshape transport texture and local class behavior, but they do not magically produce a robust braid family.

Second, sequencing and inconsistency are not automatic. The branch supports reproducible mismatch accumulation and protected ordering under external forcing, and a weaker bounded-path survivor map can open explicit kernel inconsistency. But that opening is partial and does not scale into a broad irreversibility claim.

Third, harmonic addressing behaves like a selector, not a locker. Exact-match sectors preserve braid on the clustered seed, resonance alignment stratifies into discrete address bands, and a narrow detuning threshold appears near $\delta \approx 0.75$. But bounded latches do not convert that selection into a locked or quasi-bounded exchange regime.

The final C0 boundary can therefore be stated without hype:
\begin{quote}
\emph{The no-distance combinatorial line supports graph-structural class survival, limited protected sequencing, and weak harmonic-address topology selection, but it does not yet support broad irreversibility scaling or topology locking.}
\end{quote}

That is not a failure of the line. It is exactly the kind of bounded map a control program is supposed to produce.

\end{document}
